\section{Scope, method, and conceptual framework}
\subsection{A constructive study of life-like processes}
Artificial life asks what kinds of organization can arise when components are
given explicit rules of interaction, reproduction, development, or maintenance.
The field's computational (soft), robotic (hard), and biochemical (wet) approaches
are overlapping experimental methods, rather than successive levels of
achievement \citep{Aguilar2014}. A robot can be evolved in simulation and then
manufactured; a biochemical population can be controlled by software. The
substrate therefore does not, by itself, identify what a study demonstrates.

This review emphasizes soft artificial life, the setting of the TIDE project,
while retaining substantive comparisons with physical machines and chemical
systems. Its organizing questions are: what persists, what is reproduced, what
can vary, which consequences affect reproduction, and who or what supplies the
conditions for continuation? These questions make apparently similar claims
comparable without requiring a single all-or-nothing definition of life.

\subsection{Search strategy and evidence limits}
The search was conducted on 15 November 2026. It began with cached studies of
Tierra, Avida, evoloops, Stringmol, computational replicators, and open-ended
evolution, then followed three parallel streams: recent frontiers, stagnation
and measurement, and spatial/update semantics. Additional searches covered
cellular-automata foundations, embodied systems, and chemical construction.
Original papers, author and institutional copies, and official repositories
were preferred. Field reviews provide taxonomy; experimental claims are tied
to the relevant primary work. Exact queries and reading-depth notes accompany
the LaTeX source in the repository's research records.

This is a narrative review, not a systematic review or meta-analysis. Selection
favored studies that expose mechanisms, causal controls, or informative failure
cases. It does not establish that omitted systems lack relevant capabilities.
Versioned preprints are distinguished from journal articles. Where only
indexed primary-text passages or publisher previews were accessible, that
limitation is stated. No reported experiment was independently reproduced.
Comparisons across substrates are analytical interpretations, not quantitative
rankings on a common benchmark.

\subsection{Separate capabilities before comparing systems}
Self-organization denotes the formation of macroscopic structure through local
interactions. Self-maintenance concerns persistence under disturbance and
turnover. Self-reproduction additionally requires a causal process producing
another instance of the relevant organization. Evolution is heritable population
change, which can involve selection or drift. Natural selection specifically
concerns trait-dependent differences in reproductive success. None of these terms
automatically establishes continuing innovation \citep{Taylor2016}.

The distinctions are operational. A moving pattern can persist without
reproducing. A designed loop can reproduce without generating an evolving
population. A trained local rule can repair a shape while its parameters remain
fixed. A genetically simplified cell can divide while depending on a rich
external environment. These are informative achievements, but different ones.

Three layers should also remain distinct: the transition rules of a substrate,
the processes occurring within it, and the observer or search procedure used to
find and evaluate those processes. This distinction becomes especially useful
when comparing digital ecology with modern optimization-assisted discovery.
An external objective can be scientifically valuable without being mistaken
for a reproductive advantage generated inside the studied world.

\begin{table}[htbp]
\centering\small
\caption{Substrate classes and the questions their evidence can address. The
last column identifies a limit of inference, not a judgment of research value.}
\begin{tabularx}{\textwidth}{@{}>{\raggedright\arraybackslash}p{17mm}YYY@{}}
\toprule
Approach & Experimental object & Characteristic strength & Inference requiring more evidence \\
\midrule
Soft & Programs, automata, simulated agents and reactions & Explicit interventions, complete state access, controlled replay & Biological feasibility; physical autonomy; indefinite innovation \\
Hard & Embodied robots, manufactured structures and collectives & Contact with material dynamics, sensing, actuation and fabrication & Autonomous reproduction or evolution beyond an external optimizer \\
Wet & Compartments, catalysts, genomes and cells & Actual chemical compatibility, material turnover and biochemical heredity & A complete autonomous life cycle from separately demonstrated modules \\
\bottomrule
\end{tabularx}
\label{tab:substrates}
\end{table}

